% Part: lambda-calculus
% Chapter: syntax
% Section: terms

\documentclass[../../../include/open-logic-section]{subfiles}

\begin{document}

\olfileid[hi]{lam}{syn}{trm}
\olsection{पद}

लैम्ब्डा कलन के पद चरों $\Obj{v_0}$, $\Obj{v_1}$, \dots{} के अनंत भंडार,
चिह्न~``$\lambd$'' और कोष्ठकों से आगमनतः बनाए जाते हैं। चरों को दर्शाने के
लिए हम $x$, $y$, $z$, \dots{} और पदों को दर्शाने के लिए $M$, $N$, $P$,
\dots{} उपयोग करेंगे।

\begin{defn}[पद] \ollabel{defn:term}
लैम्ब्डा कलन के \emph{पदों} का समुच्चय आगमनतः इस प्रकार परिभाषित होता है:
\begin{enumerate}
  \item \ollabel{defn:term-var} यदि $x$ कोई चर है, तो $x$ पद है।
  \item \ollabel{defn:term-abs} यदि $x$ कोई चर और $M$ कोई पद है, तो
    $(\lambd[x][M])$ पद है।
  \item \ollabel{defn:term-app} यदि $M$ और $N$ दोनों पद हैं, तो $(MN)$ पद
    है।
\end{enumerate}
\end{defn}

यदि पद $(\lambd[x][M])$ को \olref{defn:term-abs} के अनुसार बनाया गया है,
तो हम उसे किसी \emph{अमूर्तन} का परिणाम कहते हैं और $\lambd[x]$ में $x$ को
\emph{!!{parameter}} कहते हैं। \olref{defn:term-app} के अनुसार बनाया गया
पद $(MN)$ किसी \emph{अनुप्रयोग} का परिणाम है।

ऊपर परिभाषित पदों में सभी कोष्ठक लिखे गए हैं। यह काफी बोझिल हो सकता है,
जैसा पद $(\lambd[x][((\lambd[x][x])(\lambd[x][(xx)]))])$ दिखाता है। हम
कोष्ठकों से बचने के कुछ नियम प्रस्तुत करेंगे। फिर भी, औपचारिक परिभाषा से
यह निर्धारित करना सरल है कि \olref{defn:term} के अनुसार कोई पद कैसे बना
है। उदाहरणार्थ, पद $(\lambd[x][((\lambd[x][x])(\lambd[x][(xx)]))])$ को
बनाने का अंतिम चरण ऐसा अमूर्तन होना चाहिए जिसका !!{parameter}~$x$ है। यह
पद $((\lambd[x][x])(\lambd[x][(xx)]))$ से अमूर्तन द्वारा मिलता है, और वह
दो पदों का अनुप्रयोग है। उन दोनों पदों में से प्रत्येक स्वयं किसी अमूर्तन
का परिणाम है, और इसी प्रकार आगे।

\begin{prob}
$(\lambd[g][(\lambd[x][(g (x x))])
  (\lambd[x][(g (x x))])])$ का निर्माण वर्णित करें।
\end{prob}

\end{document}
